\documentclass[11pt]{article}

\usepackage[a4paper,margin=30mm]{geometry}
\usepackage{amsmath,amssymb,amsthm}
\usepackage[T1]{fontenc}
\usepackage{microtype}
\usepackage[hidelinks]{hyperref}

\newtheorem{theorem}{Theorem}
\newcommand{\F}{\mathbb F}

\title{One Isomorphism Class among Six Deleted-Slope\\
Neumaier Graphs of Order 280}
\author{Xiaojian Zeng\\
\small PBC School of Finance, Tsinghua University, Beijing 100190, China\\
\small \texttt{xiaojianzeng@foxmail.com}}
\date{}\hypersetup{pdftitle={One Isomorphism Class among Six Deleted-Slope Neumaier Graphs of Order 280},pdfauthor={Xiaojian Zeng}}

\begin{document}
\maketitle

\begin{abstract}
Greaves and Tan construct a Neumaier graph from a Hadamard matrix and a
family of mutually appropriate Latin squares.  Fix the normalized Paley
Hadamard matrix of order eight and the six Desarguesian squares of order
seven, and form six graphs by deleting one square at a time.  We prove that
these six graphs form exactly one graph-isomorphism class.  If the deleted
slopes are $a,b\in\F_7^*$ and $u=b/a$, an explicit isomorphism is
$(\alpha,r,x)\mapsto(u\alpha,u^{-1}r,x)$.  The proof is accompanied by exact
adjacency lists, all 36 ordered vertex permutations, and a dependency-free
verifier.  The result is confined to this fixed six-deletion problem; it does
not resolve the broader dependence of the construction on arbitrary input
matrices and Latin-square families.
\end{abstract}

\section{The fixed classification problem}

Greaves and Tan define $\Gamma(H;\mathbf L)$ from a normalized Hadamard
matrix $H$ and a family $\mathbf L$ of mutually appropriate Latin squares
(MALS) \cite{GreavesTan2026}.  Their row--symbol interchange converts the
standard Desarguesian mutually orthogonal Latin squares into the MALS
\begin{equation}\label{eq:latin}
 L_\alpha(r,c)=\alpha(r-c),
 \qquad \alpha\in\F_7^*,\quad r,c\in\F_7.
\end{equation}
For any five of these six squares, their Theorem~4.3 gives a Neumaier graph
with parameters
\[
 (v,k,\lambda;e,s)=(280,160,96;18,35)
\]
and five distinct adjacency eigenvalues.  Their Example~5.1 records this
smallest parameter set, while their Question~5.4 asks more broadly how the
isomorphism type depends on the Hadamard matrix and the MALS family
\cite{GreavesTan2026}.  Common parameters and spectrum do not classify the
six deletions considered here.

We now give the complete labeled construction used in this note.  Write
$X=\{\infty\}\cup\F_7$ and let $\chi$ be the quadratic character of $\F_7$,
with $\chi(0)=0$.  The Paley construction \cite{Paley1933} gives the
normalized order-eight Hadamard matrix whose first row is all $1$ and whose
row indexed by $t\in\F_7$ is
\[
 h_t(\infty)=1,
 \qquad
 h_t(c)=
 \begin{cases}
 -1,&c=t,\\
 \chi(c-t),&c\ne t.
 \end{cases}
\]
Put $B_t=h_t^{\mathsf T}h_t$.  For a deleted slope $a\in\F_7^*$, set
\[
 S_a=\F_7^*\setminus\{a\},
 \qquad
 \Omega_a=S_a\times\F_7\times X.
\]
When $\alpha\ne\beta$, the rows $r$ of $L_\alpha$ and $s$ of $L_\beta$
agree in a unique column.  Their common symbol is
\begin{equation}\label{eq:K}
 K_{\alpha\beta}(r,s)
   =\frac{\alpha\beta(r-s)}{\beta-\alpha}\in\F_7.
\end{equation}
The graph $G_a$ has vertex set $\Omega_a$.  Two distinct vertices
$(\alpha,r,x)$ and $(\beta,s,y)$ are adjacent precisely when
\begin{equation}\label{eq:adjacency}
 \begin{cases}
 r\ne s, & \alpha=\beta,\\
 B_{K_{\alpha\beta}(r,s)}(x,y)=1, & \alpha\ne\beta.
 \end{cases}
\end{equation}
This is the specialization of the $A_4$ and $A_2$ relations in the source
construction \cite{GreavesTan2026}.

The problem is to classify $G_a$ for $a\in\F_7^*$.  It is a fixed finite
specialization of Question~5.4 in \cite{GreavesTan2026}; no assertion about
the priority or status of other specializations is made here.

\section{The isomorphisms}

\begin{theorem}\label{thm:classification}
The six graphs $G_a$, $a\in\F_7^*$, form exactly one isomorphism class.  For
every ordered pair $a,b\in\F_7^*$, put $u=ba^{-1}$.  Then
\begin{equation}\label{eq:phi}
 \phi_{a,b}:\Omega_a\longrightarrow\Omega_b,
 \qquad
 (\alpha,r,x)\longmapsto(u\alpha,u^{-1}r,x)
\end{equation}
is a graph isomorphism.
\end{theorem}

\begin{proof}
Multiplication by $u$ maps the unique omitted slope $a$ to $b$, and hence
bijectionally maps $S_a$ onto $S_b$.  Multiplication by $u^{-1}$ permutes
$\F_7$, while the coordinate in $X$ is fixed.  Thus \eqref{eq:phi} is a
bijection.

Suppose first that two vertices have the same slope.  Slope equality and row
inequality are both preserved by \eqref{eq:phi}, so the first relation in
\eqref{eq:adjacency} is preserved.  If $\alpha\ne\beta$, direct substitution
in \eqref{eq:K} gives
\begin{align*}
 K_{u\alpha,u\beta}(u^{-1}r,u^{-1}s)
 &=\frac{(u\alpha)(u\beta)u^{-1}(r-s)}{u(\beta-\alpha)}\\
 &=\frac{\alpha\beta(r-s)}{\beta-\alpha}
 =K_{\alpha\beta}(r,s).
\end{align*}
Consequently, before and after applying $\phi_{a,b}$, the second relation in
\eqref{eq:adjacency} consults the same matrix $B_t$ at the same pair $(x,y)$.
It too is preserved.  The inverse coordinate map is $\phi_{b,a}$, so
$\phi_{a,b}$ is an isomorphism.  Since $a$ and $b$ were arbitrary, all six
graphs are isomorphic.  The family is nonempty, and therefore the number of
isomorphism classes is exactly one.
\end{proof}

The calculation notably does not require a permutation of the Hadamard rows:
the common symbol in \eqref{eq:K} is unchanged, rather than merely relabeled.
It therefore works for the fixed field-symbol-to-row assignment above.  The
same calculation would work for any other assignment held fixed across the
six graphs, but Theorem~\ref{thm:classification} makes no comparison between
different assignments or different Hadamard matrices.

\section{Representative and exact certificates}

We designate $G_1$ with the displayed Paley matrix convention as the
canonical representative of the single class.  Its retained slopes are
ordered as $2,3,4,5,6$; rows use the order $0,1,\ldots,6$; and the last
coordinate uses the order $\infty,0,1,\ldots,6$.  If
$\operatorname{pos}(\alpha)\in\{0,\ldots,4\}$ is the zero-based position of
$\alpha$ in the retained-slope list and $\operatorname{ind}(x)$ is the
zero-based position of $x$ in the last-coordinate list, the vertex identifier
is
\[
 8\bigl(7\operatorname{pos}(\alpha)+r\bigr)
   +\operatorname{ind}(x)\in\{0,\ldots,279\}.
\]
Equation~\eqref{eq:adjacency}, together with this ordering, determines its
adjacency list without a graph-canonicalization convention.

The ancillary files \path{evidence/adjacency/G_1.adj} through
\path{G_6.adj} contain this list and the analogous lists for all six
deletions.  The file
\path{evidence/slope_scaling_certificates.json} serializes each of the
36 ordered maps $\phi_{a,b}$ as a 280-entry target-vertex permutation, and
\path{evidence/verification-summary.json} records the deterministic check.
The accompanying Python verifier has no third-party dependencies.  For each deletion it
constructs the graph in two ways: first by scanning all seven columns for the
unique equality in \eqref{eq:latin}, and second by using \eqref{eq:K}.  It
requires equality of the two full edge sets, simplicity, symmetric adjacency,
and constant degree.  It then verifies bijectivity and equality of image
neighbor sets for every serialized permutation.  A complete run reports, for
each graph, 280 vertices, degree 160, and 22,400 edges, and accepts all 36
ordered certificates.  All arithmetic is exact modulo seven; neither
floating-point arithmetic nor randomization is used.

This certificate pattern---serialized finite objects together with a short
program that reconstructs their defining data---also appears in exact finite
graph computation by Zeng \cite{Zeng2026SSRN7385138}.  In the present note the
computation checks the construction and the published permutations; it is
corroboration rather than a substitute for the symbolic proof of
Theorem~\ref{thm:classification}.

From the project root, the deterministic verification command is
\begin{verbatim}
python3 evidence/verify_slope_scaling.py --verify-existing
\end{verbatim}
The command recomputes the graphs and certificates in memory and requires the
saved adjacency lists, certificate collection, and summary to match byte for
byte.

\section{Scope and assistance disclosure}

The conclusion classifies exactly the six graphs obtained by deleting one
member of the complete order-seven Desarguesian family while keeping the
order-eight Paley matrix fixed.  It does not answer when unrelated Hadamard
matrices, unrelated MALS families, or other orders yield isomorphic graphs.

OpenAI Codex assisted with the proof development, implementation of the exact
verifier, and preparation of the manuscript.  A separate model context
reconstructed the symbolic argument and ran an independently written
in-memory check of the adjacency lists and permutations.  That process is a
reproducibility audit, not external peer review.  The theorem rests on the
explicit field calculation above; the computational claims are limited to
the exhaustive exact checks described in the preceding section.

\enlargethispage{2\baselineskip}
\bibliographystyle{alpha}
\bibliography{references}

\end{document}
